% Quelle: Klausurvorbereitung 1 - 8 Lösungen.pdf
% Art: Klausur/Klausurvorbereitung
% Themen: ISBN-10, Prüfziffer, Zahlendreher, Euklidischer Algorithmus, ggT, Lemma von Bezout, lineare diophantische Gleichung, Matrix Q, Primfaktorzerlegung, Eulersche Phi-Funktion, Turnierplan, Siebformel, Inklusion-Exklusion, RSA, Satz von Euler, Schnelles Potenzieren, simultane Kongruenzen, Chinesischer Restsatz
% Punktzahlen: in der Quelle nicht angegeben.

\section{Klausurvorbereitung 1--8 (mit Lösungen)}
%------------------------------------------------------------
\subsection{Aufgabe 1}

\begin{aufgabe}[Aufgabe 1]
Ergänzen Sie die folgende Nummer zu einer ISBN-10 Buchnummer.\\
3-05-501517-$c_{10}$.
\end{aufgabe}

\begin{loesung}
3-05-501517-$c_{10}$
% KORRIGIERT: Original mit "." als Malzeichen und "=" statt "\equiv" ohne (mod 11); Begründung für die negativen Gewichte ergänzt.
Es gilt $c_{10} \equiv \sum_{i=1}^{9} i \cdot c_i \pmod{11}$. Für die Gewichte $6, 7, 8, 9$ rechnen wir mit $-5, -4, -3, -2$ (denn $6 \equiv -5$, \dots, $9 \equiv -2 \pmod{11}$).
\begin{align*}
c_{10} &\equiv 1 \cdot 3 + 2 \cdot 0 + 3 \cdot 5 + 4 \cdot 5 + 5 \cdot 0 - 5 \cdot 1 - 4 \cdot 5 - 3 \cdot 1 - 2 \cdot 7\\
&\equiv 3 + 2 \cdot (0-7) + 3 \cdot (5-1) + 4 \cdot (5-5) + 5 \cdot (0-1)\\
&\equiv 3 - 3 + 1 - 5\\
&\equiv -4 \equiv 7 \pmod{11}
\end{align*}
Also ist $c_{10} = 7$.
\end{loesung}

%------------------------------------------------------------
\subsection{Aufgabe 2}

\begin{aufgabe}[Aufgabe 2]
Ein Wort $\mathbf{r}$ unterscheide sich von einem Codewort $\mathbf{c}$ in einem Zahlendreher an den Fehlerstellen $x$, $y$, also
\begin{align*}
\mathbf{c} &= \dots\, c_x \,\dots\, c_y \,\dots\\
\mathbf{r} &= \dots\, c_y \,\dots\, c_x \,\dots
\end{align*}
Zeigen Sie, dass \textbf{ISBN10} diesen Fehler erkennt, dass also $\mathbf{h}\,\mathbf{r}^T \neq 0$ in $\Zm{11}$.
\end{aufgabe}

\begin{loesung}
$\mathbf{e} = \mathbf{r} - \mathbf{c}$ hat an der Stelle $x$ den Eintrag $c_y - c_x$, an der Stelle $y$ den Eintrag $c_x - c_y$ und sonst $0$. Es ist $\mathbf{h} = (1, 2, \dots, 10)$.
% KORRIGIERT: Begründungen ergänzt (h c^T = 0; warum das Produkt nicht 0 ist).
\[
\Rightarrow \mathbf{h}\,\mathbf{r}^T = \mathbf{h}\,\mathbf{c}^T + \mathbf{h}\,\mathbf{e}^T = \mathbf{h}\,\mathbf{e}^T = x \cdot (c_y - c_x) + y \cdot (c_x - c_y) = (x - y) \cdot (c_y - c_x) \neq 0 \quad \text{in } \Zm{11}
\]
denn $\mathbf{h}\,\mathbf{c}^T = 0$ (Codewort), $x - y \not\equiv 0 \pmod{11}$ wegen $1 \le x < y \le 10$, $c_y - c_x \not\equiv 0 \pmod{11}$ (sonst gäbe es keinen Fehler), und $\Zm{11}$ ist ein Körper ($11$ ist prim), also ist das Produkt zweier Elemente $\neq 0$ wieder $\neq 0$.
\end{loesung}

%------------------------------------------------------------
\subsection{Aufgabe 3}

\begin{aufgabe}[Aufgabe 3]
\begin{enumerate}[label=\alph*)]
\item Berechnen Sie $\ggT(2406, 654)$.
\item Bestimmen Sie eine ganzzahlige Lösung der Gleichung
\[ 2406\,x - 654\,y = 24 \]
\item Bestimmen Sie sämtliche ganzzahlige Lösungen dieser Gleichung.
\end{enumerate}
\end{aufgabe}

\begin{loesung}
\textbf{a)} Euklidischer Algorithmus
\begin{align*}
a &= q_0\, b + r_1\\
b &= q_1\, r_1 + r_2\\
&\cdots\\
r_{n-1} &= q_n\, r_n \qquad \Rightarrow \quad r_n = \ggT(a, b)
\end{align*}
\begin{align*}
2406 &= 3 \cdot 654 + 444\\
654 &= 1 \cdot 444 + 210\\
444 &= 2 \cdot 210 + 24\\
210 &= 8 \cdot 24 + 18\\
24 &= 1 \cdot 18 + 6\\
18 &= 3 \cdot 6 \quad \Rightarrow \quad r_n = \ggT(a, b) = 6
\end{align*}

\textbf{b)} Bestimmen Sie eine ganzzahlige Lösung der Gleichung
\[ 2406\,x - 654\,y = 24 \]
24 ist ein Vielfaches von $\ggT(2406, 654)$, diese Gleichung ist daher lösbar.

\begin{enumerate}
\item Lösung: EA rückwärts
\begin{align*}
24 &= 444 - 2 \cdot 210\\
24 &= 444 - 2 \cdot (654 - 444)\\
24 &= 3 \cdot 444 - 2 \cdot 654\\
24 &= 3 \cdot (2406 - 3 \cdot 654) - 2 \cdot 654\\
24 &= 3 \cdot 2406 - 11 \cdot 654
\end{align*}
\item Lösung: Mit der Matrix $Q$
\begin{align*}
Q &= \begin{pmatrix} 3 & 1\\ 1 & 0\end{pmatrix}
\begin{pmatrix} 1 & 1\\ 1 & 0\end{pmatrix}
\begin{pmatrix} 2 & 1\\ 1 & 0\end{pmatrix}
\begin{pmatrix} 8 & 1\\ 1 & 0\end{pmatrix}
\begin{pmatrix} 1 & 1\\ 1 & 0\end{pmatrix}
\begin{pmatrix} 3 & 1\\ 1 & 0\end{pmatrix}\\
&= \begin{pmatrix} 4 & 3\\ 1 & 1\end{pmatrix}
\begin{pmatrix} 17 & 2\\ 8 & 1\end{pmatrix}
\begin{pmatrix} 4 & 1\\ 3 & 1\end{pmatrix}
= \begin{pmatrix} 4 & 3\\ 1 & 1\end{pmatrix}
\begin{pmatrix} 74 & 19\\ 35 & 9\end{pmatrix}
= \begin{pmatrix} 401 & 103\\ 109 & 28\end{pmatrix}
\end{align*}
\begin{align*}
% KORRIGIERT: Begründung für den Übergang zur nächsten Zeile ergänzt.
\operatorname{Det} Q = 401 \cdot 28 - 109 \cdot 103 &= 1 && | \cdot 6 \quad \text{(es ist } 6 \cdot 401 = 2406,\ 6 \cdot 109 = 654\text{)}\\
2406 \cdot 28 - 654 \cdot 103 &= 6 && | \cdot 4\\
2406 \cdot 112 - 654 \cdot 412 &= 24
\end{align*}
\end{enumerate}

% KORRIGIERT: Original: "Für alle weiteren Lösungen dieser Gleichung muss gelten 2406x - 654y = 0" (gilt nur für die Differenz zweier Lösungen).
\textbf{c)} Ist $(x, y)$ eine weitere Lösung, so muss für die Differenz zur Lösung $(3, 11)$ gelten
\[ 2406 \cdot (x - 3) - 654 \cdot (y - 11) = 0 \]
Es ist
\[ 2406 \cdot \frac{654}{6} - 654 \cdot \frac{2406}{6} = 0 \]
also
\[ 2406 \cdot 109 - 654 \cdot 401 = 0 \]
% KORRIGIERT: Begründung der Vollständigkeit und "t \in \Z" ergänzt.
Da $\ggT(401, 109) = 1$ ist, sind dies bis auf ganzzahlige Vielfache alle Lösungen von $2406\,x - 654\,y = 0$.
Damit erhält man für die allgemeine Lösung von $2406\,x - 654\,y = 24$
\[ x = 3 + 109\,t, \qquad y = 11 + 401\,t, \qquad t \in \Z \]
\end{loesung}

%------------------------------------------------------------
\subsection{Aufgabe 4}

\begin{aufgabe}[Aufgabe 4]
\begin{enumerate}[label=\alph*)]
\item Zerlegen Sie die Zahl $z = 360$ in Primfaktoren.
\item Bestimmen Sie $\varphi(z)$.
\end{enumerate}
\end{aufgabe}

\begin{loesung}
\begin{enumerate}[label=\alph*)]
\item $360 = 2^3 \cdot 3^2 \cdot 5$
\item $\varphi(360) = 360 \left(1 - \frac{1}{2}\right)\left(1 - \frac{1}{3}\right)\left(1 - \frac{1}{5}\right) = 360 \cdot \frac{1}{2} \cdot \frac{2}{3} \cdot \frac{4}{5} = 12 \cdot 8 = 96$\\
oder so\\
% KORRIGIERT: Voraussetzung für die Multiplikativität ergänzt.
$\varphi(360) = \varphi(8) \cdot \varphi(9) \cdot \varphi(5) = 4 \cdot 6 \cdot 4 = 96$ \quad (da $8$, $9$, $5$ paarweise teilerfremd sind)
\end{enumerate}
\end{loesung}

%------------------------------------------------------------
\subsection{Aufgabe 5}

\begin{aufgabe}[Aufgabe 5]
Bestimmen Sie einen Turnierplan für $2m = 8$ Mannschaften an 7 Spieltagen.
\end{aufgabe}

\begin{loesung}
% KORRIGIERT: Original: "x + y = 2m für x \neq y, andernfalls spielt x gegen 2m" (passt nicht zur Tabelle).
In Runde $r$ spielt $x$ gegen $y$, wenn $x + y \equiv r \pmod{2m-1}$ und $x \neq y$ ist ($x, y \in \{1, \dots, 2m-1\}$); gilt $2x \equiv r \pmod{2m-1}$, so spielt $x$ gegen $2m$.

\medskip
\begin{tabular}{ccccccc}
\toprule
Runde 1 & Runde 2 & Runde 3 & Runde 4 & Runde 5 & Runde 6 & Runde 7\\
\midrule
1\ 7 & 1\ 8 & 1\ 2 & 1\ 3 & 1\ 4 & 1\ 5 & 1\ 6\\
2\ 6 & 2\ 7 & 3\ 7 & 2\ 8 & 2\ 3 & 2\ 4 & 2\ 5\\
3\ 5 & 3\ 6 & 4\ 6 & 4\ 7 & 5\ 7 & 3\ 8 & 3\ 4\\
4\ 8 & 4\ 5 & 5\ 8 & 5\ 6 & 6\ 8 & 6\ 7 & 7\ 8\\
\bottomrule
\end{tabular}
\end{loesung}

%------------------------------------------------------------
\subsection{Aufgabe 6}

\begin{aufgabe}[Aufgabe 6]
In einem Sportverein trainieren 55 Athleten.\\
35 spielen Fußball,\\
27 Leichtathletik,\\
12 Judo,\\
13 Fußball und Leichtathletik,\\
7 Fußball und Judo,\\
5 Leichtathletik und Judo,\\
2 Fußball, Leichtathletik und Judo.\\
Wie viele Athleten trainieren eine andere Sportart?
\end{aufgabe}

\begin{loesung}
% KORRIGIERT: Klammern um die Vereinigung fehlten.
Gesucht: $\left|\Omega \setminus (A_1 \cup A_2 \cup A_3)\right|$
\begin{align*}
&= |\Omega| - \bigl(|A_1| + |A_2| + |A_3|\bigr) + |A_1 \cap A_2| + |A_1 \cap A_3| + |A_2 \cap A_3| - |A_1 \cap A_2 \cap A_3|\\
&= 55 - (35 + 27 + 12) + 13 + 7 + 5 - 2 = 80 - 76\\
&= 4
\end{align*}
\end{loesung}

%------------------------------------------------------------
\subsection{Aufgabe 7}

\begin{aufgabe}[Aufgabe 7]
RSA Verfahren. Wie berechnet man den privaten Schlüssel $d$?\\
Berechnen Sie $d$ aus der Vorlesung. In Kapitel 3 waren $m = 101$, $e = 13$.
\end{aufgabe}

\begin{loesung}
\textbf{1. Lösungsweg:} Satz von Euler, Schnelles Potenzieren
% KORRIGIERT: Original: "d = e^{\varphi(\varphi(m))-1} \equiv 1 (mod \varphi(m))" und "d = e^{39} \equiv 1 (mod 100)" (richtig: e \cdot d \equiv 1 bzw. d \equiv e^{39}); Voraussetzungen ggT = 1 und (mod 100) bei den Potenzen ergänzt.
\begin{align*}
& e \cdot d \equiv 1 \pmod{\varphi(m)}, \quad \varphi(101) = 100 \qquad (\text{Euler: } a^{\varphi(m)} \equiv 1 \pmod m \text{ für } \ggT(a, m) = 1)\\
\Rightarrow\quad & d \equiv e^{\varphi(\varphi(m)) - 1} \pmod{\varphi(m)}, \text{ denn } e \cdot d \equiv e^{\varphi(\varphi(m))} \equiv 1 \pmod{\varphi(m)}\\ & \qquad (\text{wegen } \ggT(13, 100) = 1)\\
& \varphi(\varphi(m)) = \varphi(100) = 40\\
\Rightarrow\quad & \varphi(\varphi(m)) - 1 = 39 = 1 + 2 + 4 + 32\\
& 13,\quad 13^2 \equiv 69,\quad 13^4 \equiv 69^2 \equiv 61,\quad \dots,\quad 13^{32} \equiv 81 \pmod{100}\\
\Rightarrow\quad & d \equiv e^{39} = 13^{1} \cdot 13^{2} \cdot 13^{4} \cdot 13^{32} \pmod{100}\\
\Rightarrow\quad & d \equiv 13 \cdot 69 \cdot 61 \cdot 81 \equiv 77 \pmod{100}
\end{align*}
% Hinweis: In der Quelle steht "13^2", "13^4", "69^2", "13^32" in Textschreibweise und "*" als Malzeichen.

\textbf{2. Lösungsweg:} Lemma von Bezout, EA
\[
% KORRIGIERT: Quantor "für ein y \in \Z" ergänzt.
e \cdot d - \varphi(m) \cdot y = 1 \ \text{ für ein } y \in \Z \quad\Leftrightarrow\quad e \cdot d \equiv 1 \pmod{\varphi(m)}
\]
EA mit $\varphi(m)$ und $e$
\begin{align*}
100 &= 7 \cdot 13 + 9\\
13 &= 1 \cdot 9 + 4\\
9 &= 2 \cdot 4 + 1\\
4 &= 4 \cdot 1
\end{align*}
\[
Q = \begin{pmatrix} 7 & 1\\ 1 & 0\end{pmatrix}
\begin{pmatrix} 1 & 1\\ 1 & 0\end{pmatrix}
\begin{pmatrix} 2 & 1\\ 1 & 0\end{pmatrix}
\begin{pmatrix} 4 & 1\\ 1 & 0\end{pmatrix}
= \begin{pmatrix} 8 & 7\\ 1 & 1\end{pmatrix}
\begin{pmatrix} 9 & 2\\ 4 & 1\end{pmatrix}
= \begin{pmatrix} 100 & 23\\ 13 & 3\end{pmatrix}
\]
$\operatorname{Det} Q = 100 \cdot 3 - 13 \cdot 23 = 1$\\
% KORRIGIERT: Original "d = -23"; die Determinante liefert 13 \cdot (-23) \equiv 1 (mod 100).
$\Rightarrow\quad 13 \cdot (-23) \equiv 1 \pmod{100} \quad\Rightarrow\quad d \equiv -23 \equiv 77 \pmod{100}$
\end{loesung}

%------------------------------------------------------------
\subsection{Aufgabe 8}

\begin{aufgabe}[Aufgabe 8]
Geben Sie sämtliche Lösungen der simultanen Kongruenzen an
\begin{align*}
z &\equiv 5 \pmod{13}\\
z &\equiv 4 \pmod{15}
\end{align*}
Was ist die kleinste positive Lösung?
\end{aufgabe}

\begin{loesung}
\begin{align*}
& z = m\,x + a = n\,y + b\\
% KORRIGIERT: Original "Beout" (Tippfehler); Bedeutung von m, n, a, b ergänzt.
\Rightarrow\quad & m\,x - n\,y = b - a \qquad \text{Lemma von Bezout}
\end{align*}
mit $m = 13$, $a = 5$, $n = 15$, $b = 4$ und $x, y \in \Z$.

Euklidischer Algorithmus
\begin{align*}
& 15 = 1 \cdot 13 + 2\\
& 13 = 6 \cdot 2 + 1\\
\Rightarrow\quad & 1 = 13 - 6 \cdot 2 = 13 - 6 \cdot (15 - 1 \cdot 13) = 7 \cdot 13 - 6 \cdot 15\\
\Rightarrow\quad & b - a = 4 - 5 = -1 = 6 \cdot 15 - 7 \cdot 13\\
% KORRIGIERT: Original "z = -86 (mod 195)"; richtig ist "\equiv", mit Begründung für den Modul 195.
\Rightarrow\quad & z = 5 - 13 \cdot 7 = 4 - 6 \cdot 15 = -86\\
\Rightarrow\quad & z \equiv -86 \pmod{195} \quad \text{(da } \ggT(13, 15) = 1 \text{ ist, ist } 13 \cdot 15 = 195 \text{ der Modul)}
\end{align*}
Alle Lösungen sind also $z = -86 + 195\,k$ mit $k \in \Z$. Also ist die kleinste positive Lösung
\[ \Rightarrow\quad z = -86 + 195 = 109 \]
\end{loesung}
